% META: {"id": "TSPE-GEOMETRIE-ESPACE-COURS-03", "chapitre": "TSPE-GEOMETRIE-ESPACE", "type_objet": "cours", "capacites_codes": ["C3"], "sources_inspiration": [], "mode_creation": "ex_nihilo", "status": "approved", "capacites": ["TSPE-GEOMETRIE-ESPACE-C3"], "statut": "structure"}

\section{Produit scalaire dans l'espace}

% ============================================================
% STRATE 1 — Definition, formule analytique, applications
% ============================================================

\definition[1]{
Dans un repere orthonorme $(O;\vec{i},\vec{j},\vec{k})$, si $\vec{u}\begin{pmatrix}x\\y\\z\end{pmatrix}$ et $\vec{v}\begin{pmatrix}x'\\y'\\z'\end{pmatrix}$, le \textbf{produit scalaire} de $\vec{u}$ et $\vec{v}$ est :
\[
  \vec{u}\cdot\vec{v} = xx'+yy'+zz'.
\]
}

\propriete{
$\vec{u}\cdot\vec{v} = \|\vec{u}\|\times\|\vec{v}\|\times\cos(\widehat{\vec{u},\vec{v}})$. En particulier, $\vec{u}$ et $\vec{v}$ sont orthogonaux si et seulement si $\vec{u}\cdot\vec{v}=0$.
}

\exemple{
Soit $A(1,0,2)$, $B(3,1,1)$, $C(2,-1,4)$. Montrer que le triangle $ABC$ est rectangle en $A$.

$\vec{AB}\begin{pmatrix}2\\1\\-1\end{pmatrix}$, $\vec{AC}\begin{pmatrix}1\\-1\\2\end{pmatrix}$.

$\vec{AB}\cdot\vec{AC} = 2\times1+1\times(-1)+(-1)\times2 = 2-1-2 = -1 
eq 0$.

Le triangle n'est donc \textbf{pas} rectangle en $A$.

% BEGIN-VERIFY
% AB = (2, 1, -1)
% AC = (1, -1, 2)
% dot = sum(a*b for a, b in zip(AB, AC))
% assert dot == -1
% END-VERIFY
}

\propriete[Longueur d'un vecteur]{
$\|\vec{u}\| = \sqrt{\vec{u}\cdot\vec{u}} = \sqrt{x^2+y^2+z^2}$.
}

% ============================================================
% STRATE 2 — Erreurs frequentes
% ============================================================

\erreurFrequente{
\textbf{Utiliser la formule analytique dans un repere non orthonorme.} La formule $\vec{u}\cdot\vec{v}=xx'+yy'+zz'$ n'est valable que dans un repere \textbf{orthonorme}. Toujours verifier cette hypothese avant de l'appliquer.
}
